\subsection{Fyzika a gravitace}

Unity nabízí vestavěné komponenty pro upravování gravitace a fyziky objektů. Pro základ a začátky jsou sice dostačující, ovšem pro realistický a uživatelsky přívětivý systém je potřeba vzít tyto základní komponenty, vynulovat jejich hodnoty a udělat vlastní skripty. To nabízí možnost dynamicky upravovat gravitaci podle aktuálních podmínek ve hře. Pokud chce hráč praktikovat \textit{Wall jump}, tak se gravitace trochu sníží, aby hráč nespadl příliš rychle. Zároveň lze znatelně více upravovat fyziku skoku. Ve hrách je ideální, aby trochu déle trvalo, než se hráč dostane do maximální výšky skoku. Poté je dobré, aby hráč chvíli zůstal \textit{viset} ve vzduchu a až potom rychle spadl. Za předpokladu použití vestavěných komponentů v Unity by se tato práce zbytečně znáročnila.